
\Inicio{ARTÍCULO DE INVESTIGACIÓN}{Uso de hojas de árbol para la creación de macetas biodegradables}{Zamora Perez Lesly María, Cota Velarde Xochitl Noemí, Gonzalez Sandoval Blanca Nelly}

\Resumen{Este proyecto analiza la problemática de la acumulación de hojas en zonas urbanas, como Mazatlán, donde estos residuos orgánicos pueden provocar obstrucción de drenajes, inundaciones, malos olores y afectaciones a la salud. Además, se considera el impacto ambiental del uso excesivo de macetas de plástico, las cuales generan contaminación por su lenta degradación.\\Como alternativa, se propone la elaboración de macetas biodegradables a partir de hojas secas reutilizadas, aplicando el enfoque de economía circular. La metodología consiste en triturar las hojas, mezclarlas con agua y harina como aglutinante, moldearlas y secarlas al sol para obtener un producto funcional.\\Diversos antecedentes respaldan la viabilidad de esta propuesta, demostrando que los materiales vegetales pueden sustituir al plástico sin afectar el crecimiento de las plantas. Incluso, los análisis indican que las macetas biodegradables favorecen el desarrollo radicular y aumentan la supervivencia tras el trasplante.\\En conclusión, este proyecto ofrece una solución sustentable que reduce residuos orgánicos y plásticos, contribuye al cuidado del medio ambiente y promueve prácticas ecológicas en la comunidad.}{Hojas secas, macetas biodegradables, sustentabilidad, residuos orgánicos, reducción de plásticos.}


    \begin{multicols}{2}
    

\section{Introducción}

En zonas urbanas con abundante vegetación, como Mazatlán, es frecuente la acumulación de hojas de árboles, especialmente durante ciertas temporadas del año. Aunque se trata de residuos naturales, su acumulación puede generar problemas como la obstrucción de alcantarillas, inundaciones, malos olores y afectaciones a la salud pública.

Ante esta problemática, surge la necesidad de replantear el manejo de estos residuos, considerándolos como un recurso potencial. En este sentido, el presente estudio propone la utilización de hojas secas para la elaboración de macetas biodegradables, como una alternativa ecológica que permita reducir tanto la acumulación de residuos orgánicos como el uso de plásticos.

Esta propuesta se enmarca dentro del enfoque de economía circular, promoviendo el aprovechamiento de recursos naturales y la implementación de soluciones sostenibles en el entorno urbano.

\section{Materiales y métodos}

\subsection{Materiales}

\begin{itemize}
    \item Hojas secas
    \item Bolsa de nailon
    \item 200 g de harina
    \item Vinagre
    \item Agua 
    \item Moldes para macetas
    \item Recipiente
\end{itemize}

\subsection{Procedimiento}

Las hojas secas fueron trituradas hasta obtener una textura fina. Posteriormente, se añadió agua y harina como aglutinante, formando una mezcla homogénea. Esta mezcla se vertió en moldes con la forma deseada. Finalmente, las macetas fueron expuestas a la luz solar hasta su completo secado y endurecimiento.

\section{Resultados}

Se evaluaron dos grupos experimentales: 6 macetas de plástico (MP) y 6 macetas biodegradables (MB). Las macetas biodegradables mostraron un mejor desempeño en el crecimiento inicial de las plantas en comparación con las macetas de plástico (\autoref{tab:resultados_dos_grupos}).

    \begin{table}[H]

\caption{Resultados de los dos grupos}
\label{tab:resultados_dos_grupos}

\begin{threeparttable}

{\centering\small\setstretch{1.35}\rowcolors{2}{Grey}{White}
\begin{tabular}{M{0.3\linewidth}M{0.275\linewidth}M{0.275\linewidth}}\hline

    \theader{Variable} & \theader{Plástico (MP)} & \theader{Biodegradable (MB)} \\ \hline
    
    Número de plantas emergidas & 2 de 6       & 3 de 6 \\
    Altura de la planta (cm)    & 1 y 2        & 2.5, 2.8 y 3 \\
    Altura promedio (cm)        & 1.5          & 2.77 \\
    Raíz seca                   & 20\% y 40\%  & 10\%, 20\% y 20\% \\
    Promedio de raíz seca       & 30\%         & 16.7\% \\
    Raíces deformadas           & 0\%          & 0\% \\
    
    \hline
\end{tabular}}

\begin{tablenotes}\tiny % ex: \tnote{1} ... \item[1]
    \item[] 
\end{tablenotes}

\end{threeparttable}
\end{table}

No se observaron deformaciones radiculares en ninguno de los tratamientos. El análisis estadístico indicó que las diferencias no fueron significativas (p > 0.05), probablemente debido al tamaño reducido de la muestra.

\section{Discusión}

Los resultados sugieren que las macetas biodegradables elaboradas a partir de hojas secas pueden proporcionar condiciones favorables para el crecimiento vegetal, particularmente en el desarrollo radicular. Aunque no se alcanzó significancia estadística, las tendencias observadas indican ventajas en términos de retención de humedad y aireación.

Estos hallazgos coinciden con estudios previos que destacan el potencial de los materiales vegetales como sustitutos del plástico en aplicaciones agrícolas. Además, la reutilización de hojas contribuye a la reducción de problemas ambientales urbanos, como la obstrucción de drenajes y la acumulación de residuos.
No obstante, se recomienda realizar estudios con muestras más amplias para confirmar estos resultados.

\section{Conclusión}

Las macetas biodegradables elaboradas a partir de hojas secas representan una alternativa viable frente a las macetas de plástico. A pesar de que los resultados no fueron estadísticamente significativos, se observó una tendencia favorable en el crecimiento vegetal y la calidad radicular.

Esta propuesta contribuye a la reducción de residuos orgánicos y plásticos, promoviendo prácticas sostenibles y el aprovechamiento de recursos naturales dentro de un enfoque de economía circular.


% ------------------------------------------------------ %

\bibliographystyle{apalike-ejor}
\bibliography{art_04/bib}

\end{multicols}